\chapter{Shaders: Four Answers to XNA's .fx}
\label{ch:shaders-four-answers}
\label{ch:shader-fx-gap}
\index{shader}\index{.fx}

There is no single CNA answer to XNA effects.  The current tree contains four different
shader-delivery strategies: CNA-owned implementations of the stock semantics, recompilation
of Microsoft's stock HLSL sources on D3D9, execution of Microsoft's compiled stock bytecode
through FNA3D and MojoShader, and the public renderer-specific \cnaclass{ShaderEffect} route.
This chapter separates those strategies before naming the narrower compatibility gap that
remains: a game still cannot load an arbitrary user-authored compiled \texttt{.fx} effect and
drive its reflected parameters through the common CNA renderer contract.

\begin{figure}[tbp]
\centering
\begin{figurealt}{Four distinct shader routes. XNA-shaped stock effects produce GpuDrawParams consumed by renderer-owned implementations. FNA3D internally executes retained stock bytecode through FNA3D and MojoShader. ShaderEffect accepts renderer-specific source or program forms. CNA graphics extensions provide separate post-processing effects. The public arbitrary Effect bytecode constructor is an explicit refusal.}
\resizebox{\textwidth}{!}{%
\begin{tikzpicture}[node distance=9mm and 13mm]
  \node[cna public] (stock) {five stock effects};
  \node[cna internal, right=of stock] (params) {\texttt{GpuDrawParams}};
  \node[cna internal, right=of params] (owned) {renderer-owned stock path};
  \node[cna internal, below=of stock] (blobs) {retained stock \texttt{.fxb} blobs};
  \node[cna internal, right=of blobs] (fna3d) {FNA3D + MojoShader\\internal stock route};
  \node[cna public, below=of blobs] (shader) {\texttt{ShaderEffect}};
  \node[cna internal, right=of shader] (program) {renderer-specific source,\\bytecode, or program};
  \node[cna public, below=of shader] (ext) {CNA graphics extensions};
  \node[cna internal, right=of ext] (post) {post-process/material routes};
  \node[cna warning, above=of owned] (refuse) {\texttt{Effect(device, byte[])}\\throws};
  \draw[cna flow] (stock) -- (params);
  \draw[cna flow] (params) -- (owned);
  \draw[cna flow] (blobs) -- (fna3d);
  \draw[cna flow] (shader) -- (program);
  \draw[cna flow] (ext) -- (post);
  \draw[cna optional] (refuse) -- node[cna label]{not a portable route} (owned);
\end{tikzpicture}}
\end{figurealt}
\caption{CNA's shader and effect strategies are separate APIs and evidence scopes. FNA3D's
internal stock-bytecode path does not implement the public arbitrary-bytecode constructor.}
\label{fig:shader-routes}
\end{figure}



\section{XNA's compiled-effect format}

A normal XNA game does not ship \texttt{.fx} shader source text. The XNA content pipeline
compiles \texttt{.fx} effect files to a proprietary compiled bytecode format at build time,
and the compiled bytecode --- not the HLSL source --- is what ships inside the game's
\texttt{.xnb} content files and gets loaded at runtime via
\texttt{new Effect(GraphicsDevice, byte[])}. On the real Xbox 360/Windows XNA runtime, that
bytecode is interpreted directly; FNA instead ships MojoShader, a from-scratch bytecode
\emph{parser and cross-compiler} that turns the same compiled Direct3D~9 shader bytecode into
GLSL (or other target shading languages) at load time, which is exactly what lets FNA run
historical XNA effect content without its original source.

\section{The public custom-effect boundary}

Two public entry points state the boundary directly.  The \cnaclass{Effect} constructor taking
a byte buffer ignores that buffer and always throws \texttt{NotImplementedException}
(Chapter~\ref{ch:stock-effects}); the content registry recognizes the general XNA
\texttt{EffectReader} but installs a known-unsupported reader that throws
\texttt{ContentLoadException} with the reason ``compiled platform shader bytecode.''  Stock
effect readers are registered separately and do work.

The missing piece is larger than parsing bytecode.  CNA's common draw boundary is the closed
\cnaclass{GpuDrawParams} aggregate: each built-in effect fills named C\texttt{++} fields and each
renderer switch-selects one of its already known programs.  It has no reflected arbitrary
parameter table, technique/pass state, or common parameter-to-uniform binding protocol.  A
general parser would therefore still need a portable representation and binding path for the
program it decoded.  The project's planning documents
(\texttt{docs/shader-effect-vs-fx-bytecode.md} and \texttt{docs/fx-bytecode-support-plan.md})
scope that broader work, but their historical dependency claims must not be mistaken for the
current source tree: MojoShader is already present for the deliberately closed FNA3D stock path
described below.

Do not attach a current sample count to that gap without retesting the samples.  The pinned
\texttt{cna-samples} catalog contains 153 source-archive entries split among completed ports,
tracked placeholders, and deliberate exclusions (Chapter~\ref{ch:samples-and-examples}); its
dated \texttt{missing.md} files are case records, not fresh proof that every old blocker remains.
For a concrete port, inventory each general \texttt{EffectReader} asset and direct bytecode
constructor call instead.

\section{Answer one: reimplement the stock semantics}

Most CNA renderers do not consume XNA bytecode.  The five public stock effects translate their
properties into \cnaclass{GpuDrawParams}; each renderer then chooses CNA-authored shader variants
for combinations such as fog, vertex color, alpha test, and lighting.  This is broad practical
coverage, but it is a semantic reimplementation rather than execution of Microsoft's compiled
programs.  It also explains why adding a new arbitrary parameter to an effect is not merely a
map insertion: the aggregate, every relevant renderer, and the shader variants must agree on it.

\section{Answer two: recompile Microsoft's sources on D3D9}

The Direct3D~9 renderer takes a different evidentiary route.  It vendors the six Microsoft/FNA
stock \texttt{.fx} sources --- the five public effects plus the internal sprite effect --- and
their four \texttt{.fxh} includes, discovers entry points from the source's own
\texttt{compile vs\_2\_0}/\texttt{ps\_2\_0} statements, and compiles them with
\texttt{d3dcompiler\_47}.  Its comparison tool strips only bytecode comment blocks such as CTAB
and the compiler creator string: 61 of 66 resulting shader instruction streams match the
Microsoft-shipped programs exactly.  The five misses are all pixel-lighting vertex variants;
the recorded investigation attributes them to the compiler-47 versus XNA-era compiler-43
difference after testing the available optimization flags.  This is unusually strong stock
effect evidence, but it still does not compile a game's arbitrary \texttt{.fx} content.

\section{Answer three: execute the stock bytecode through FNA3D}

FNA3D is the important counterexample to any claim that CNA contains no bytecode compatibility.
The repository commits six FNA-derived XNA~4.0 \texttt{.fxb} files: Sprite, Basic, AlphaTest,
DualTexture, EnvironmentMap, and Skinned effects, 106~KB in total.  The build embeds those exact
blobs; renderer initialization calls \texttt{FNA3D\_CreateEffect}, which runs them through the
pinned MojoShader dependency and exposes their parameter metadata.  Draw dispatch selects the
original program's shader variant and applies its named parameters.

This route is intentionally closed.  FNA3D exposes no source-string compiler, CNA's FNA3D
\texttt{CreateEffectRenderer()} returns null, and the renderer reports
\texttt{GraphicsCapability::CustomEffects == false}.  The committed blobs are also runtime
artifacts rather than in-repository build products because producing \texttt{.fxb} requires the
old \texttt{fxc} toolchain.  FNA3D therefore proves that real XNA stock bytecode executes in CNA;
it does not make the public arbitrary-effect constructor or general XNB reader work.

\section{Answer four: bypass \texttt{.fx} with ShaderEffect}

CNA does not leave custom shading entirely unaddressed, but ``cross-renderer'' needs a precise
definition here. \cnaclass{ShaderEffect} (\texttt{CNAEXT}) presents one constructor shape and
one \cnaclass{IEffectRenderer} contract to every renderer. A renderer factory override proves
only that an effect object is returned. It does \emph{not} prove that arbitrary source is
compiled, that every setter has meaning, or that a 3D draw consumes the resulting program.
Reading the live implementations and their registered tests gives this narrower matrix:

\begin{longtable}{@{}p{0.18\textwidth}p{0.76\textwidth}@{}}
\toprule
\textbf{Renderer} & \textbf{Input and path that actually consumes it} \\
\midrule
\endhead
\textbf{EasyGL} &
GLSL source is compiled and linked by the active GL driver. \cnaclass{SpriteBatch} and the
general 3D draw paths consume it. Uniform names and array setters are real; 2D, cube, and
volume-texture binds each have a dedicated pixel test. \\
\textbf{SDL GPU} &
GLSL source is compiled at runtime to SPIR-V through \texttt{libshaderc}, then consumed by
\cnaclass{SpriteBatch}'s fixed 32-byte vertex contract and 128-byte uniform block. This includes
a tested two-output MRT shader. \\
\textbf{Vulkan} &
The two \texttt{std::string} arguments carry raw, precompiled SPIR-V bytes --- not GLSL source.
\cnaclass{SpriteBatch} consumes them through a fixed 32-byte vertex contract and 128-byte
push-constant layout, verified by \texttt{Vulkan\_ShaderEffect\_SpirV}. \\
\textbf{D3D9} &
HLSL is compiled at runtime to the active SM2/SM3 profile and consumed by
\cnaclass{SpriteBatch}, with a real color-inversion pixel test. \\
\textbf{D3D11/12} &
HLSL is compiled at runtime as \texttt{vs\_5\_0}/\texttt{ps\_5\_0} and consumed by
\cnaclass{SpriteBatch}'s fixed vertex and 128-byte constant-buffer contracts. Scalar/vector
setter names are ignored; the sprite texture and sampler still arrive through
\texttt{t0}/\texttt{s0}. \\
\textbf{BGFX} &
No input is accepted through this API: \texttt{CompileProgram()} always returns false. The
object reports that BGFX requires precompiled binary shaders, but \cnaclass{ShaderEffect}
exposes no route that loads those binaries. \\
\textbf{SDL renderer} &
There is no effect renderer, so construction leaves \texttt{IsEffectValid()} false. Passing the
non-null object to \cnaclass{SpriteBatch}\texttt{::Begin()} throws
\texttt{std::runtime\_error}; merely constructing it does not throw. \\
\bottomrule
\end{longtable}

WebGPU and FREEDIRECT/free-direct likewise inherit the null \texttt{CreateEffectRenderer()} default
rather than compiling a public custom program; their \cnaclass{SpriteBatch} implementations
reject a non-null custom effect. Canvas rejects it for the same no-programmable-stage reason.
The Software and Headless overrides serve different testing purposes, not real shader
execution: Software accepts any non-empty pair and marks it valid while continuing to render
through its fixed CPU shading path, whereas Headless records compilation, bind, uniform, and
2D-texture activity without producing shader pixels. These are useful contracts, but calling
either one custom-shader support would be misleading.

None of these \cnaclass{ShaderEffect} paths is bytecode compatibility. A game cannot reuse an existing compiled
\texttt{.fx} blob: it must provide renderer-specific source for EasyGL, SDL GPU, and Direct3D,
or precompile the exact SPIR-V contract Vulkan expects. Even among the working implementations,
the public setter names and supported draw shapes are not portable promises. The API is thus a
real, pixel-verified alternative on a bounded set of paths, not a write-once shader format for
all CNA renderers.

\cnaclass{GpuDrawParams} (Chapter~\ref{ch:stock-effects}) does carry a
\texttt{customEffectRenderer} pointer. \cnaclass{ShaderEffect} populates it through its
\texttt{FillGpuDrawParams()} override. At present, however, only EasyGL's general 3D draw
dispatch reads that field.
Vulkan and SDL GPU retain similarly named custom-effect pointers inside their
\emph{SpriteBatch-specific} snapshots; the Direct3D custom paths are also SpriteBatch
facilities. The default-no-op setters on \cnaclass{IEffectRenderer}
(Chapter~\ref{ch:backend-architecture}) are another deliberate capability boundary: a public
\texttt{SetUniformXxx()} call may be structurally accepted without changing GPU state unless
that concrete renderer overrides it.

\subsection{Object lifetime, submission, and cloning are deliberately narrow}

The constructor takes its two source arguments by reference but immediately copies both into
the effect. It is therefore safe to construct from temporary or subsequently modified strings;
the program is compiled from the construction-time copies, not from caller-owned character
storage. If the device has a renderer, construction asks it once for an
\cnaclass{IEffectRenderer}. A returned but invalid renderer does not make the common constructor
throw: it writes that renderer's \texttt{GetCompileError()} text to standard error and leaves
\texttt{IsEffectValid()} false. A null renderer and a failed compile are intentionally collapsed
to the same public boolean, and there is no public \cnaclass{ShaderEffect} accessor for the
diagnostic. Renderers may of course fail before returning from their own factory, so this is not a
promise that every renderer-specific setup failure becomes a boolean.

Treat that validity result as a precondition, not a delayed exception. \cnaclass{Effect}\texttt{::Apply()}
does reject an already disposed effect, but its common code calls \texttt{OnApply()} and then
makes the effect current. \cnaclass{ShaderEffect}\texttt{::OnApply()} merely writes a debug
message when its renderer is absent or invalid; it does not bind a fallback program. Calling
\texttt{Apply()} on an invalid custom effect can therefore still replace the device's current
effect pointer. There is no automatic retry or source recompilation in this object; construct a
new effect after correcting source or changing the renderer circumstances.

The uniform and texture calls are submission calls, not an XNA-style parameter collection.
\cnaclass{ShaderEffect} keeps no uniform dictionary, no sampler table, and no texture ownership:
each \texttt{SetUniformXxx()} forwards its name and scalar or caller buffer immediately to the
renderer if one exists, while each \texttt{SetTexture()} immediately forwards the address of the
texture's renderer. It performs no common-layer check of a uniform name, sampler-unit range,
array count, or matrix/array pointer. EasyGL uploads/binds at that call; some other renderers copy
into a fixed staging block; inherited interface defaults can do nothing; and Headless records a
raw texture-renderer pointer for tracing. Thus a successful call is not a portable guarantee of a
binding, and it never gives the effect ownership of the texture. Keep every texture alive by its
normal owner and set it again whenever the target renderer/path requires rebinding. The standalone
setters also do not inspect \texttt{IsDisposed}; unlike \texttt{Apply()}, they continue to
forward while the renderer object remains alive, so disposal must be treated as terminal by the
caller.

\texttt{Clone()} is equally source-only. It returns an owning raw \cnaclass{Effect} pointer
whose concrete object is constructed afresh as
\texttt{new ShaderEffect(*device\_, vertSrc\_, fragSrc\_)}: it recompiles an independent
renderer program and restores the three \cnaclass{IEffectMatrices} values to identity. It does
not copy submitted uniforms, texture binds, name/tag metadata, or a compiled-program handle;
its base effect begins with its own empty parameter collection and one Default/P0 technique.
Use an owning smart pointer at the call site and reapply all desired state to the clone. The
existing \texttt{ShaderEffectTests} test establishes distinct object identity, copied source
strings, and equal validity on the default test device. It does not compile against a renderer or
cover matrix/state transfer, texture lifetime, invalid-apply behavior, or a renderer reset.

\subsection{A minimal GLSL ShaderEffect adapted from a test}

Adapted from \texttt{easygl\_shadereffect\_texturecube\_test.cpp}, a minimal
cube-map-sampling fragment shader and its construction:

\begin{lstlisting}
const char* vertSrc = R"(#version 300 es
precision highp float;
layout(location = 0) in vec3 aPosition;
void main() { gl_Position = vec4(aPosition, 1.0); }
)";

const char* fragSrc = R"(#version 300 es
precision highp float;
out vec4 FragColor;
uniform samplerCube CubeSampler;
uniform vec3 direction;
void main() { FragColor = texture(CubeSampler, direction); }
)";

ShaderEffect effect(getGraphicsDeviceProperty(), vertSrc, fragSrc);
if (!effect.IsEffectValid()) {
    // compilation failed -- check the renderer's own log output before drawing
}

effect.SetUniformInt("CubeSampler", 0);
effect.SetUniformVec3("direction", dx, dy, dz);
effect.SetTexture(0, environmentCube); // TextureCube overload
\end{lstlisting}

Check \texttt{IsEffectValid()} before drawing. Unlike a stock effect selected from CNA's shipped
programs or FNA3D's retained stock bytecode, \cnaclass{ShaderEffect} compiles game-supplied source
at runtime. A syntax error or unsupported GLSL feature therefore produces an invalid effect that
must be reported before submission.

\subsection{The D3D9 HLSL counterpart}

\cnaclass{ShaderEffect}'s constructor takes one pair of source strings for every renderer; there
is no HLSL-specific overload. On a Direct3D renderer the same call shape compiles HLSL instead of
GLSL. The following color-inversion example is adapted from the pinned D3D9 test and uses its
SM2/SM3 \texttt{Macros.fxh} texture macros, which the vendored stock
\texttt{SpriteEffect.fx} also uses:

\begin{lstlisting}
const char* vertexSrc = R"(
float2 vpSize;
struct VSInput { float3 Position : POSITION0; float4 Color : COLOR0; float2 TexCoord : TEXCOORD0; };
struct VSOutput { float4 Position : POSITION0; float4 Color : COLOR0; float2 TexCoord : TEXCOORD0; };
VSOutput main(VSInput input) {
    VSOutput output;
    float2 ndc = (input.Position.xy / vpSize) * 2.0 - 1.0;
    output.Position = float4(ndc.x, -ndc.y, 0.0, 1.0);
    output.Color = input.Color;
    output.TexCoord = input.TexCoord;
    return output;
}
)";

const char* pixelSrc = R"(
texture2D Texture;
sampler TextureSampler : register(s0) = sampler_state { Texture = (Texture); };
float4 main(float4 color : COLOR0, float2 texCoord : TEXCOORD0) : COLOR0 {
    float4 texColor = tex2D(TextureSampler, texCoord);
    return float4(1.0 - texColor.rgb, 1.0); // the same inversion the GLSL example performs
}
)";

ShaderEffect invertEffect(getGraphicsDeviceProperty(), vertexSrc, pixelSrc);
// ... same IsEffectValid() check as the GLSL example ...

SamplerState pointClamp = SamplerState::PointClamp;
spriteBatch_->Begin(SpriteSortMode::Deferred, BlendState::Opaque,
                     &pointClamp, nullptr, nullptr, &invertEffect);
spriteBatch_->Draw(texture, destRect, Color::White);
spriteBatch_->End();
\end{lstlisting}

\texttt{vpSize} is worth pausing on, because it is doing real work rather than being decorative:
a custom vertex shader supplied this way has no other way to map \cnaclass{SpriteBatch}'s
pixel-space \texttt{Position} input into normalized device coordinates, so this uniform (set
automatically by the renderer, not by game code) is what makes the sprite land at the correct
screen location at all --- the same real test this example is adapted from separately verifies
that sampling \emph{outside} the sprite's destination rectangle stays the clear color, proving
this NDC math is genuinely correct rather than happening to paint something plausible-looking
over the whole screen. One small, harmless documentation inaccuracy worth knowing if you go
reading \cnaclass{ShaderEffect}'s own header next: its class-level doc comment describes it as
a ``GLSL-source-based effect.'' That is accurate for EasyGL and SDL GPU, but this exact worked
example proves the same class compiles HLSL on Direct3D, while Vulkan interprets the strings as
raw SPIR-V bytes. The comment is stale cross-renderer documentation, not a constraint enforced
by the public constructor.

\subsection{Beyond SpriteBatch: a 3D draw call}

Both worked examples so far apply a \cnaclass{ShaderEffect} through
\cnaclass{SpriteBatch}. The wider path described in this subsection is specifically
\textbf{EasyGL support}; no other renderer's general 3D draw dispatch currently reads
\texttt{GpuDrawParams::customEffectRenderer}. For a long stretch of this project's history that was the only
place a \cnaclass{ShaderEffect} could actually take effect: applying one before a raw
\cnaclass{GraphicsDevice}\allowbreak\texttt{::DrawIndexedPrimitives()} call --- the call
shape \cnaclass{ModelMesh}\allowbreak\texttt{::Draw()} and a hand-rolled
\cnaclass{RawMesh} both use --- was silently ignored. Because
\cnaclass{ShaderEffect} did not yet implement \cnaclass{IEffectMatrices}
(Chapter~\ref{ch:stock-effects}) and its \texttt{FillGpuDrawParams()} override still ran
the \cnaclass{Effect} base class's no-op, \texttt{GpuDrawParams} stayed at its default,
zero value on every 3D draw, and \textbf{EasyGL}'s
\texttt{DrawIndexedPrimitivesEx()} fell back to selecting one of its own built-in,
stride-keyed shaders regardless of whatever custom effect the caller had just applied ---
exactly the same silent-fallback failure shape Chapter~\ref{ch:graphicsdevice}'s
\texttt{ExtractMatrices()} section already documents for stock effects, one layer further
out.

This was closed by three coordinated changes, verified directly against
\texttt{ShaderEffect.hpp} / \texttt{.cpp} and \texttt{EasyGLRenderer.cpp}: (1)
\texttt{GpuDrawParams} gained a \texttt{customEffectRenderer} field; (2)
\cnaclass{ShaderEffect} now implements \cnaclass{IEffectMatrices} --- real
\texttt{World} / \texttt{View} / \texttt{Projection} properties, extracted by the same
\texttt{GraphicsDevice::ExtractMatrices()} every stock effect already goes through --- and
overrides \texttt{FillGpuDrawParams()} to populate \texttt{customEffectRenderer}; (3)
EasyGL's \texttt{DrawPrimitivesEx()} / \texttt{DrawIndexedPrimitivesEx()} check that field
first and, when it is set, bind the custom compiled program and its
\texttt{World} / \texttt{View} / \texttt{Projection} uniforms directly, bypassing the
built-in stride-selected shaders entirely.

A worked example, adapted from \texttt{easygl\_shadereffect\_3d\_test.cpp}, makes the
wiring concrete. The shader is a minimal N-dot-L diffuse lighting pass over a single
textured quad, using the same GLSL attribute-location convention (\texttt{location = 0/1/2}
for \texttt{Position} / \texttt{Normal} / \texttt{TexCoord}) as
\cnaclass{VertexPositionNormalTexture}'s existing 32-byte stride:

\begin{lstlisting}
auto* fx = dynamic_cast<ShaderEffect*>(fxBase_.get());
fx->setWorldProperty(world);
fx->setViewProperty(Matrix::CreateLookAt(Vector3(0, 0, 3), Vector3::Zero, Vector3::Up));
fx->setProjectionProperty(Matrix::CreatePerspectiveFieldOfView(
    MathHelper::PiOver4, aspectRatio, 0.1f, 100.0f));
fx->Apply();                          // binds the compiled program first
fx->SetTexture(0, whiteTex_);
fx->SetUniformVec3("uLightDir", 0.0f, 0.0f, 1.0f);
fx->SetUniformVec3("uDiffuseColor", 200.0f / 255.0f, 100.0f / 255.0f, 50.0f / 255.0f);

device.SetVertexBuffer(vb_.get());
device.setIndicesProperty(ib_.get());
device.DrawIndexedPrimitives(PrimitiveType::TriangleList, 0, 0, 4, 0, 2);
\end{lstlisting}

The real test behind this example is deliberately built to rule out a false positive: it
draws the same quad twice, once with \texttt{World = Identity} (facing the light, world
normal \texttt{(0,0,1)}, expected color close to the full \texttt{(200,100,50)} diffuse
value) and once with \texttt{World = CreateRotationY(180\textdegree{})}. A
\(90\)\textdegree{} rotation was deliberately \emph{not} chosen for the second case ---
it would turn the quad edge-on to the camera, so ``renders nothing'' and ``renders
black'' would look identical in a single-pixel readback. \(180\)\textdegree{} keeps the
exact same on-screen footprint (the quad is X-symmetric, so mirroring X does not change
its silhouette) while flipping the world-space normal to \texttt{(0,0,-1)}, which the
shader's own \texttt{max(dot(...), 0.0)} clamps to a genuinely lit black. Only a
\texttt{World} matrix that is actually reaching the vertex shader and actually affecting
world-space lighting can produce that second, specific result rather than merely happening
to draw something.

\subsection{A second, narrower gap this one uncovered, then closed: arbitrary layouts}

The first version of the 3D draw path above recognized only the five byte-strides EasyGL's
then-existing \texttt{ApplyLayout()} switch knew about (16/20/24/32/52). Later PBR and
vertex-color work expanded the fixed family to 48/56/68 too, but a stride is still not a
description of its internal fields. Arbitrary layouts were deliberately scoped out of the
first task, since several real ported sample shaders
(\texttt{NormalMapping.fx}'s \texttt{Position+Normal+Binormal+Tangent+TexCoord} layout,
among others) need element arrangements that match none of those fixed interpretations. A
truly custom \cnaclass{VertexDeclaration} --- arbitrary element count, arbitrary per-element
offset, any \cnaclass{VertexElementFormat} --- reads back correctly through the same 3D
\cnaclass{ShaderEffect} draw path via a second, separately-tracked fix:
\texttt{VertexBuffer::SetDataRaw()}
now pushes the buffer's own \cnaclass{VertexDeclaration} down to the renderer through a new
\texttt{IVertexBuffer}\allowbreak{}\texttt{Renderer::}\allowbreak{}%
\texttt{SetVertexDeclaration()} call (default no-op, so
Vulkan/BGFX/SDL\_Renderer needed zero code changes for this to stay correct), and
\texttt{ApplyLayout()} binds
generically from that declaration when one is present --- GLSL attribute location equals
the element's own index within the declaration --- falling straight through to the
original stride-keyed switch, unchanged, whenever no declaration was supplied. This is a
purely additive capability, not a replacement of the fixed-stride path the first worked
example already exercises.

The real test behind this (\path{easygl_shadereffect_custom_vertex_layout_test.cpp})
is worth citing for its own verification technique, not just its result: it declares a
5-element, 48-byte layout (\texttt{Position} / \texttt{Normal} / \texttt{Tangent} / %
\texttt{TextureCoordinate} / \texttt{Color}). CNA now also has a fixed 48-byte PBR layout,
but that one contains a four-float tangent and no trailing color; sharing a total stride does
not make the records equivalent. The test diagnostically encodes three of the four
non-position attributes directly into the output pixel
(\texttt{FragColor = vec4(Normal.x, Tangent.y, TexCoord.x, Color.r)}) and draws the same
quad twice with two independently-chosen sets of attribute values. Distinct, correctly-offset
readback values across both draws is what proves each attribute is being read from its own
correct byte offset --- not aliased to a neighbor, and not a value that happens to be
correct by coincidence.

\section{The honest summary}

If your porting target relies only on \cnaclass{SpriteBatch} and the five public stock effects
from Chapter~\ref{ch:stock-effects}, it can use CNA's renderer-owned implementations; FNA3D even
executes the six original stock blobs, including the internal sprite program.  If the port loads
a game-authored compiled \texttt{.fx} shader, the public path still stops at the throwing
constructor or known-unsupported content reader on every renderer.  What changes per renderer is
whether a \cnaclass{ShaderEffect} rewrite is usable and which contract it must target: EasyGL
takes GLSL and alone extends the effect to general 3D draws; SDL GPU takes GLSL through shaderc;
Vulkan takes precompiled SPIR-V; D3D9/11/12 take HLSL for SpriteBatch; BGFX has no valid public
path; and the non-programmable or not-yet-wired renderers reject, simulate, or only record the
object.  The accurate statement is therefore neither ``CNA has no XNA bytecode support'' nor
``effects are compatible.''  Stock semantics, stock source recompilation, closed stock-bytecode
execution, and renderer-specific custom programs are real; arbitrary compiled effects plus
general parameter plumbing remain the bounded gap.
